\chapter{Conclusioni} \label{chap:conclusioni}

\textit{In questo capitolo si presenta il consuntivo finale del lavoro svolto,
valutando il grado di raggiungimento degli obiettivi e di soddisfazione dei
requisiti di progetto, e si propone, oltre alle conoscenze acquisite, una 
riflessione personale sul percorso intrapreso.}

\section{Consuntivo finale}
\label{sec:consuntivo-finale}

La Tabella \ref{tab:consuntivo-ore} riporta la ripartizione delle 320 ore di stage tra le fasi in cui il
lavoro è stato organizzato, descritte nella \hyperref[sec:pianificazione]{Sezione Pianificazione del Capitolo}.

\begin{table}[H]
    \centering
    \begin{tabular}{p{9cm} c}
        \hline
        \textbf{Attività principali} & \textbf{Ore} \\
        \hline
        Onboarding e analisi del progetto & 16 \\
        \hline
        Studio delle tecnologie e analisi del firmware & 56\\
        \hline
        Analisi e progettazione delle modifiche & 64\\
        \hline
        Sviluppo e ottimizzazione dello storage & 24\\
        \hline
        Gestione dei dati e comunicazioni & 40\\
        \hline
        Analisi tecnica CAN sniffer e integrazione firmware & 56\\
        \hline
        Testing e validazione & 48\\
        \hline
        Documentazione finale & 16\\
        \hline
        \multicolumn{1}{r}{\textbf{Totale}} & \textbf{320} \\
        \hline
    \end{tabular}
    \caption{Ripartizione delle ore di stage per fase}
    \label{tab:consuntivo-ore}
\end{table}

\section{Raggiungimento degli obiettivi}
\label{sec:raggiungimento-obiettivi}

Questa sezione mette a confronto gli obiettivi definiti dall'azienda in fase di pianificazione (Sezione 
\hyperref[chap:descrizione-stage]{Obiettivi}) con quanto effettivamente realizzato.

\subsection{Obiettivi obbligatori}

\begin{center}
\begin{longtable}{|p{1.6cm}|p{11cm}|p{2.2cm}|}
\hline
\multicolumn{1}{|c|}{\textbf{Codice}} & \multicolumn{1}{c|}{\textbf{Obiettivo}} & \multicolumn{1}{c|}{\textbf{Raggiunto}} \\
\hline
\endfirsthead
\rowcolor{white}
\multicolumn{3}{c}{{\bfseries \tablename\ \thetable{} - Continua dalla pagina precedente}}\\
\hline
\multicolumn{1}{|c|}{\textbf{Codice}} & \multicolumn{1}{c|}{\textbf{Obiettivo}} & \multicolumn{1}{c|}{\textbf{Raggiunto}} \\
\hline
\endhead
\hline
\rowcolor{white}
\multicolumn{3}{|r|}{{Continua nella prossima pagina...}}\\
\hline
\endfoot
\endlastfoot

\textbf{OO-1} & \textit{Refactoring} e ottimizzazione del sistema di lettura e scrittura su scheda SD & Sì \\
\hline
\textbf{OO-2} & Completamento del sistema di rotazione dei file mediante l'utilizzo di ring-buffer & Sì \\
\hline
\textbf{OO-3} & Implementare la scrittura e salvataggio dei dati raccolti direttamente in formato binario & Sì \\
\hline
\textbf{OO-4} & Implementazione di un sistema di \textit{recovery} che permette di riprendere la scrittura su SD in caso di guasti o rimozione della stessa dal suo \textit{slot} & Sì \\
\hline
\textbf{OO-5} & Miglioramento dei servizi BLE e introduzione di nuovi comandi firmware & Sì \\
\hline
\textbf{OO-6} & Aggiunta di metriche di confronto per poter stimare il miglioramento rispetto alla versione precedente del firmware & Sì \\
\hline
\textbf{OO-7} & Implementare un protocollo di calibrazione dei sensori & Sì \\
\hline
\textbf{OO-8} & Aggiungere una logica di sessionamento, attivando i sensori solamente all'avvio della sessione e disattivandoli quando termina la sessione & Sì \\
\hline
\textbf{OO-9} & Produzione della documentazione tecnica relativa al firmware sviluppato & Sì \\
\hline

\caption{Raggiungimento degli obiettivi obbligatori}
\label{tab:obiettivi-obbligatori}
\end{longtable}
\end{center}

\subsection{Obiettivi desiderabili}

\begin{table}[H]
    \centering
    \begin{tabular}{|p{1.6cm}|p{11cm}|p{2.2cm}|}
        \hline
        \textbf{Codice} & \textbf{Obiettivo} & \textbf{Raggiunto} \\
        \hline
        \textbf{OD-1} & Permettere la visualizzazione dei movimenti della moto in tempo reale con un modellino 3D & Sì \\
        \hline
        \textbf{OD-2} & Integrazione dell'analisi e gestione dei dati provenienti dalla ECU & Sì \\
        \hline
        \textbf{OD-3} & Implementare un sistema di rilevamento di assenza della scheda microSD & Sì \\
        \hline
    \end{tabular}
    \caption{Raggiungimento degli obiettivi desiderabili}
    \label{tab:obiettivi-desiderabili}
\end{table}

\subsection{Obiettivi facoltativi}

\begin{table}[H]
    \centering
    \begin{tabular}{|p{1.6cm}|p{11cm}|p{2.2cm}|}
        \hline
        \textbf{Codice} & \textbf{Obiettivo} & \textbf{Raggiunto} \\
        \hline
        \textbf{OF-1} & Implementazione di un sistema di indicizzazione delle sessioni memorizzate sulla scheda SD & Sì \\
        \hline
        \textbf{OF-2} & Implementazione di un sistema di \textit{inactivity} per permettere la disattivazione dei sensori quando non viene rilevato alcun movimento per un tempo configurabile dall'utente & Sì \\
        \hline
    \end{tabular}
    \caption{Raggiungimento degli obiettivi facoltativi}
    \label{tab:obiettivi-facoltativi}
\end{table}

\subsection{Requisiti raggiunti}

Sulla base delle scelte progettuali e implementative descritte nei capitoli precedenti, risultano
raggiunti complessivamente \textbf{37} requisiti sui \textbf{38} tracciati in \hyperref[sec:tracciamento_requisiti]{Sezione Tracciamento dei requisiti}. 
La Tabella \ref{tab:requisiti-raggiunti} riporta, per ciascuna categoria e urgenza, il numero di requisiti
raggiunti sul totale tracciato.

\begin{center}
\begin{tabular}{|l|c|c|c|c|}
\hline
 & \textbf{Obbligatori} & \textbf{Desiderabili} & \textbf{Opzionali} & \textbf{TOTALE} \\
\hline
RF & 14/14 & 7/7 & 1/1 & 22/22 \\
\hline
RN & 3/4 & 3/3 & 1/1 & 7/8 \\
\hline
RV & 8/8 & 0/0 & 0/0 & 8/8 \\
\hline
\textbf{TOTALE} & \textbf{25/26} & \textbf{10/10} & \textbf{2/2} & \textbf{37/38} \\
\hline
\end{tabular}
\captionof{table}{Requisiti raggiunti, per categoria e urgenza}
\label{tab:requisiti-raggiunti}
\end{center}

L'unico requisito non pienamente raggiunto è \textbf{\hyperref[req:rn1]{R-N-Ob-1}}, che riguarda 
la riduzione del consumo di RAM rispetto alla versione precedente del firmware. La motivazione 
dietro al parziale raggiungimento è stata approfondita maggiormente nella sezione sulla 
\hyperref[subsubsec:ram-occupazione]{Occupazione della RAM}.

\section{Conoscenze acquisite}
\label{sec:conoscenze-acquisite}
Nel corso del tirocinio sono state acquisite e consolidate competenze tecniche e metodologiche. 
In particolare, è stata maturata esperienza nello sviluppo e nel refactoring di firmware embedded 
in C++ su ESP32, con gestione della concorrenza tramite FreeRTOS e variabili atomiche, nella 
progettazione di formati binari a basso overhead con meccanismi di recovery e nella realizzazione 
di servizi BLE, sia in modalità connessa sia broadcast. Sono inoltre state sviluppate competenze 
nell’analisi di protocolli CAN proprietari tramite tecniche di sniffing, nella valutazione di 
alternative hardware secondo criteri di costo e rischio e nello sviluppo mobile con React 
Native/TypeScript, con attenzione all’ottimizzazione del rendering. Le scelte implementative sul 
firmware sono state inoltre valutate tramite benchmark comparativi, misurando quantitativamente 
l’impatto su CPU e RAM. \\

Dal punto di vista metodologico, il tirocinio ha permesso di consolidare la capacità di lavorare 
su componenti differenti di uno stesso sistema, coordinandone le interfacce, di operare secondo 
una metodologia Agile e di condurre analisi preventive dei rischi. È stata inoltre sviluppata una 
maggiore capacità di motivare le scelte progettuali e di affrontare attività di debug in ambito 
embedded, ricercando le cause dei problemi anche a livello fisico ed elettrico.

\section{Valutazione personale}
\label{sec:valutazione-personale}

L'esperienza di tirocinio svolta presso Eggon è stata, nel complesso, \textbf{estremamente positiva}. 

Un aspetto che ho apprezzato particolarmente è stato lo sviluppo congiunto di firmware e
applicazione insieme alla mia collega tirocinante: questo lavoro parallelo, ma costantemente
allineato, ci ha tenuti uniti per l'intera durata del tirocinio ed è stato per me un'occasione
importante di crescita nella capacità di comunicare in modo chiaro ed efficace il proprio
lavoro, soprattutto durante gli \textit{stand-up meeting} giornalieri. \\

Mi sono inoltre trovato molto bene con il team di Eggon, che ringrazio per il supporto costante, 
per aver stimolato in me la ricerca di soluzioni sempre migliori e per la disponibilità all'ascolto 
anche nei momenti di difficoltà, specialmente durante le fasi di implementazione più complesse e 
ricche di bug.

Nel complesso, ritengo questa esperienza di tirocinio molto formativa, sia dal punto di vista
tecnico sia da quello relazionale.

\newpage